\begin{tikzpicture}[font=\figurefont\small,
  box/.style={draw=BookRule,fill=BookPaper,align=center,
    text width=31mm,minimum height=14mm},
  arr/.style={-{Stealth[length=2mm]},draw=BookInk,line width=.8pt},
  feedback/.style={arr,dashed,draw=BookTeal}]
  \node[box,fill=FigureInput!12] (task) at (0,2) {状态与外在任务\\到右房间目标G};
  \node[box,fill=FigureOutput!12] (high) at (5,2) {高层$\zeta$\\选择下门或Option};
  \node[box,fill=FigureOutput!12] (low) at (5,-.7) {低层$\pi_\omega$\\每步选择方向};
  \node[box,fill=FigureInput!12] (env) at (10,-.7) {房间网格\\执行原子动作};
  \node[box,fill=FigureModel!12] (end) at (0,-.7) {终止$\beta_\omega$\\到门后交还控制};
  \node[box,fill=FigureLoss!12] (reward) at (10,2) {原始任务奖励\\可选子目标奖励};
  \draw[arr] (task) -- (high);
  \draw[arr] (high) -- node[right] {编号或目标} (low);
  \draw[arr] (low) -- (env);
  \draw[arr] (env) -- (reward);
  \draw[feedback] (reward) -- (high);
  \draw[feedback] (reward.south west) -- (low.north east);
  \draw[arr] (env.south) -- (10,-2.3) -- (0,-2.3) -- (end.south);
  \node[fill=white] at (5,-2.3) {新状态同时供低层与终止判定使用};
  \draw[arr] (5,-2.3) -- (low.south);
  \draw[arr] (end.north) -- (0,.7) -- (3,.7) -- (high.south west);
  \node[align=center,text width=135mm] at (5,-3.5)
    {高层调用间隔可随机；低层每步仍观测状态。\\学习低层会改变高层动作的转移含义，旧经验需要相应处理。};
\end{tikzpicture}
